%==============================================================
\chapter{In-Field Rotation of CS2 and OCS}
\label{ch:results}
%==============================================================

\section{Experimental}
\subsection{Breaking the wall}
\begin{figure}[ht]
    \centering
    \includegraphics[width=\textwidth]{04_experiments/css_jet_and_droplets.png}
    \caption{}
    \label{fig:css_jet_and_droplets}
\end{figure}
\begin{figure}[ht]
    \centering
    \includegraphics[width=\textwidth]{04_experiments/ocs_jet_and_droplets.png}
    \caption{}
    \label{fig:ocs_jet_and_droplets}
\end{figure}
\begin{itemize}
    \item First experiments with css and ocs -> breaking the wall
    \item explain the jet in detail: why are the oscillations so good, what are we expecting from a jet scan, clean spectrogram
    \item make insert for the spectrogram to explain what values we use
    \item Differences in css and ocs regarding their envelope and why we have a different slower centrifuge for ocs
    \item mention the other experiments with DIB and NO-Dimer
\end{itemize}

\subsection{Decelerating the molecules}
\begin{itemize}
    \item plots 
    \begin{itemize}
        \item css plot with jet decelerating and then droplets accelerating and decelerating all with spectrogram
        \item ocs the same
    \end{itemize}
    \item talk about why we did decelerating -> because the theory developed for explaining the differences between ocs and css implied independence 
    \item say what we see qualitatively 
    \item tell why it's not a surprise that the jet scans can be decelerated
    \item 
\end{itemize}
\section{Theory from Austria}

\section{Discussion}
\label{sec:discussion}
\begin{figure}[ht]
    \centering
    \includegraphics[width=\textwidth]{04_experiments/discussion.png}
    \caption{}
    \label{fig:discussion}
\end{figure}
\begin{itemize}
    \item plots 
    \begin{itemize}
        \item css and ocs accelerating and decelerating overlaid with the theoretical data
    \end{itemize}
\end{itemize}
